\chapter{Project budget}
\label{chap:budget}

This chapter gives the direct cash budget and the resources the project used
without paying for them. The project was planned around computers the
authors already had, open-source engineering software, Autodesk educational
access and digital submission, so its direct cash budget is zero.

\section{Direct cash budget}
\label{sec:budget}

No equipment, software licence, cloud service, data, printing, travel,
fabrication or testing was paid for from project funds.
Table~\ref{tab:cash-budget} sets out the reasons.

\begin{table}[H]
\centering
\caption{Direct cash budget}
\label{tab:cash-budget}
\small
\begin{tabularx}{\textwidth}{
    >{\raggedright\arraybackslash}p{0.23\textwidth}
    >{\raggedright\arraybackslash}X
    >{\centering\arraybackslash}p{0.15\textwidth}}
\toprule
\textbf{Category} & \textbf{Why no cash is needed} & \textbf{Budget (KES)}\\
\midrule
Computers & Existing personal and university computers. The finite element models were kept small enough to solve on them, the largest at 209,763 unknowns in under a minute. & 0\\[0.35em]
CAD and analysis software & build123d, OpenCASCADE, Gmsh, scikit-fem and PyAMG are open source and provide the automated geometry and stress-analysis pipeline locally. & 0\\[0.35em]
Autodesk Fusion & Educational access, used for early exploratory trials and for the optional CAD connector. & 0\\[0.35em]
Other software & Python, NumPy, SciPy, FastAPI, ReportLab, React, three.js, Git and LaTeX are free to use. & 0\\[0.35em]
Cloud services & Everything runs on the user's computer. No paid hosting or remote computation is used. & 0\\[0.35em]
Internet and electricity & Existing personal and university services. & 0\\[0.35em]
Printing and travel & The report is submitted digitally. & 0\\[0.35em]
Fabrication and testing & Physical fabrication and load testing are outside the scope of this project. & 0\\[0.35em]
Contingency & Schedule margin and scope control instead of a cash reserve. & 0\\
\midrule
\textbf{Total} & & \textbf{0}\\
\bottomrule
\end{tabularx}
\end{table}

\section{In-kind resources}
\label{sec:in-kind-resources}

A zero cash budget does not mean the project used no resources.
Table~\ref{tab:in-kind} lists them. They are not given money values, because
there is no sound basis for pricing them.

\begin{table}[H]
\centering
\caption{In-kind resources}
\label{tab:in-kind}
\small
\begin{tabularx}{\textwidth}{
    >{\raggedright\arraybackslash}p{0.25\textwidth}
    >{\raggedright\arraybackslash}X
    >{\raggedright\arraybackslash}X}
\toprule
\textbf{Resource} & \textbf{How it was provided} & \textbf{What it was used for}\\
\midrule
Computers & Personal and university & Development, automated tests, case studies, report\\
Open-source software & Public licences & Geometry, meshing, finite element analysis, server, application, report\\
Autodesk Fusion & Educational access & Early CAD and simulation trials; optional parametric padeye connector\\
Design codes and catalogues & University library and public sources & Rule sets, material grades, shackle and pipe reference data\\
Internet and electricity & Personal and university & Research, package downloads, running the tool\\
Student time and supervision & Academic project & Marine engineering analysis, programming, verification and writing\\
\bottomrule
\end{tabularx}
\end{table}

\section{Costs beyond this project}
\label{sec:budget-beyond}

Taking a design from \NEMO{} into shipboard service requires expenditure outside
this academic project. A proof load test of a lifting padeye, vibration
measurements on an engine-room pipe run, model basin or sea trials of a
stabilizer fin, certified marine steel plate, fabrication and non-destructive
weld inspection each carry their own yard and survey costs. A commercial
organisation adopting the tool would also need licensed copies of the
referenced design codes: ASME BTH-1, AISC 360 and AWS D1.1 are sold by their
publishers, whereas IACS Unified Requirements are freely published.
